\section*{Abstract}
\addcontentsline{toc}{section}{Abstract}
\markboth{Abstract}{}

% English abstract (200--300 words):
% context, objectives, method used, and main results.

\begin{otherlanguage}{english}
Traditional data migration processes, particularly during the analysis and design phases, suffer from a heavy reliance on manual tasks and fragmented management tools (such as spreadsheets and decentralized office documents). This classic approach, prone to human error, results in extended delays and a lack of traceability that limits overall project efficiency.

The objective of this project is to design and develop the ADMP AI Studio, an intelligent platform aimed at modernizing system architecture and automating data migration ETL (Extract, Transform, Load) pipelines within the Accenture group. The goal is to drastically reduce manual intervention while standardizing the generation of transformation rules.

The adopted methodology relies on deploying an innovative Agentic Artificial Intelligence architecture. A multi-agent system, orchestrated via LangGraph and integrating advanced Large Language Models (LLMs) using "Chain of Thought" reasoning, was implemented. Technically, the architecture transitions toward centralized storage under PostgreSQL leveraging the JSON format. Various specialized agents (Scope, Design, Validation) work in synergy to ingest data models, autonomously perform Gap Analysis, and validate the consistency of operations.

The main results highlight the successful automation of critical design deliverables (Data Mapping, Table Mapping, SDS, TDS). The new infrastructure establishes a continuous, secure, and auditable data flow, thereby ensuring data quality before injection into target systems, and paving the way for the automation of complex transformation phases.
\end{otherlanguage}

\vspace{1em}
\noindent\textbf{Keywords:} Data migration, Agentic Artificial Intelligence, ETL pipeline, System architecture, Automation.
